\documentclass[11pt]{article}

\usepackage[margin=1in]{geometry}
\usepackage{amsmath,amssymb,amsthm}
\usepackage{array}
\usepackage{enumitem}
\usepackage{hyperref}

% Common PDE notation
\newcommand{\R}{\mathbb{R}}
\newcommand{\lap}{\Delta}
\newcommand{\grad}{\nabla}

\title{APM5557: Partial Differential Equations\\Weak and Strong Maximum Principles}
\author{}
\date{}

\begin{document}

\maketitle

Both principles say a harmonic function cannot exceed its boundary values.
\textbf{The weak version says the maximum is reached on the boundary, though it
might also be reached inside. The strong version says it cannot be reached
inside at all, unless $u$ is constant.}

Throughout, $u$ is harmonic ($\lap u = 0$) in a bounded open set $\Omega$ and
continuous on $\overline{\Omega}$. Applying either principle to $-u$ gives the
same result for minimums.

\section*{The two statements}

\textbf{Weak maximum principle:}
\[
\max_{\overline{\Omega}} u = \max_{\partial\Omega} u
\]

\textbf{Strong maximum principle:} if $\Omega$ is connected and $u$ reaches its
maximum at an interior point, then $u$ is constant.

\section*{Comparison}

\begin{center}
\renewcommand{\arraystretch}{1.4}
\begin{tabular}{|>{\raggedright\arraybackslash}p{0.2\textwidth}
                |>{\raggedright\arraybackslash}p{0.34\textwidth}
                |>{\raggedright\arraybackslash}p{0.34\textwidth}|}
\hline
 & \textbf{Weak} & \textbf{Strong} \\
\hline
\textbf{Claim} & The maximum occurs on $\partial\Omega$
  & The maximum occurs only on $\partial\Omega$, unless $u$ is constant \\
\hline
\textbf{Interior maximum} & Not ruled out & Forces $u$ to be constant \\
\hline
\textbf{Needs $\Omega$ connected?} & No & Yes \\
\hline
\textbf{Usual proof} & Perturbation and calculus & Mean value property \\
\hline
\textbf{Typical use} & Uniqueness, stability, comparing solutions
  & Strict inequalities such as $u > 0$ inside \\
\hline
\end{tabular}
\end{center}

The strong version implies the weak one. If the maximum is reached inside, the
component of $\Omega$ containing that point has constant $u$, so the same value
is also reached on the boundary.

\section*{Why each is true}

\textbf{Weak:} let $v = u + \varepsilon|x|^2$, so $\lap v = 2n\varepsilon > 0$.
At an interior maximum of $v$, all second derivatives would be $\le 0$, so
$\lap v \le 0$, which is a contradiction. So $v$ reaches its maximum on the
boundary, and letting $\varepsilon \to 0$ gives the result for $u$.

\textbf{Strong:} suppose $u$ reaches its maximum $M$ at an interior point. By
the mean value property, $M$ is the average of $u$ over a small ball around
that point. Since $u \le M$ everywhere, that average can only equal $M$ if
$u = M$ on the whole ball. So the set where $u = M$ is open. It is also closed
in $\Omega$ by continuity, so connectedness makes it all of $\Omega$.

\section*{Why the hypotheses matter}

\begin{itemize}
\item \textbf{Connectedness:} let $\Omega$ be two separate disks, with $u = 1$
  on one and $u = 0$ on the other. $u$ is harmonic and reaches its maximum
  inside, but is not constant. The weak principle still holds, but the strong
  one fails.
\item \textbf{Boundedness:} let $\Omega$ be the upper half-plane with $u = y$.
  Then $u = 0$ on the boundary, but $u$ grows without limit inside, so even the
  weak principle fails.
\end{itemize}

\section*{What each gives you}

\begin{itemize}
\item \textbf{Weak:} if $u_1$ and $u_2$ solve the same Dirichlet problem, then
  $w = u_1 - u_2$ has zero boundary values. Its maximum and minimum are both 0,
  so $w = 0$. This gives uniqueness.
\item \textbf{Strong:} suppose $\Omega$ is connected, $u \ge 0$ on
  $\partial\Omega$, and $u$ is not identically zero there. The weak principle
  gives only $u \ge 0$ inside, but the strong principle gives $u > 0$
  everywhere inside. If $u$ were 0 at an interior point, that would be an
  interior minimum, which would make $u$ constant.
\end{itemize}

\end{document}
